\documentclass{article}
\usepackage{amsmath, amssymb, amsthm}
\newtheorem{theorem}{Theorem}
\newtheorem{claim}{Claim}
\newtheorem{definition}{Definition}
\newtheorem{assumption}{Assumption}
\begin{document}

\begin{theorem}  % Goal
Prove that $\frac{1}{1+a^2} + \frac{1}{1+b^2} + \frac{1}{1+c^2} \le \frac{3}{2}$.
\end{theorem}

\section*{Given}
\begin{assumption}  % Assumption
$a, b, c > 0$.
\end{assumption}

\begin{assumption}  % Assumption
$a + b + c = abc$.
\end{assumption}

\section*{Derived Statements}
\begin{claim}  % Step 1
$a,b,c > 0$, $a+b+c = abc$ implies $a^2 + b^2 + c^2 \ge 9$, with equality if and only if $a = b = c = \sqrt{3}$.
\end{claim}
\begin{proof}
Let $s = a+b+c = abc$. By AM–GM,
\[
s = a+b+c \ge 3\sqrt[3]{abc} = 3\sqrt[3]{s},
\]
so $s^3 \ge 27 s$, i.e., $s^2 \ge 27$. Hence $s \ge 3\sqrt{3}$ (since $s>0$). By Cauchy–Schwarz,
\[
(a^2+b^2+c^2)(1^2+1^2+1^2) \ge (a+b+c)^2 = s^2,
\]
thus
\[
a^2+b^2+c^2 \ge \frac{s^2}{3} \ge \frac{(3\sqrt{3})^2}{3} = 9.
\]
Equality in AM–GM requires $a=b=c$; equality in Cauchy–Schwarz also requires $a=b=c$. Substituting $a=b=c$ into $a+b+c=abc$ gives $3a = a^3$, whence $a = \sqrt{3}$. Therefore equality holds iff $a=b=c=\sqrt{3}$.
\end{proof}

\begin{claim}  % Step 2
$a,b,c>0$, $a+b+c=abc$ implies $ab+bc+ca\ge9$, with equality if and only if $a=b=c=\sqrt{3}$.
\end{claim}
\begin{proof}
From $a+b+c=abc$, divide by $abc$ (positive) to get $\frac{1}{ab}+\frac{1}{bc}+\frac{1}{ca}=1$. Set $p=ab$, $q=bc$, $r=ca$. Then $p,q,r>0$ and $\frac{1}{p}+\frac{1}{q}+\frac{1}{r}=1$. By the AM–HM inequality,
\[
\frac{p+q+r}{3}\ge\frac{3}{\frac{1}{p}+\frac{1}{q}+\frac{1}{r}}=\frac{3}{1}=3,
\]
hence $p+q+r\ge9$, i.e., $ab+bc+ca\ge9$. Equality in AM–HM occurs iff $p=q=r$. Combined with $\frac{1}{p}+\frac{1}{q}+\frac{1}{r}=1$, this gives $p=q=r=3$, so $ab=bc=ca=3$. Then $ab=bc$ implies $a=c$, and $bc=ca$ implies $b=a$; thus $a=b=c$. From $ab=3$ we obtain $a=\sqrt{3}$. Conversely, if $a=b=c=\sqrt{3}$, then $ab+bc+ca=3+3+3=9$.
\end{proof}

\begin{claim}  % Step 3
$a,b,c>0$, $a+b+c=abc$ implies $\dfrac{1}{\sqrt{1+a^{2}}}+\dfrac{1}{\sqrt{1+b^{2}}}+\dfrac{1}{\sqrt{1+c^{2}}}\le\dfrac{3}{2}$, with equality if and only if $a=b=c=\sqrt{3}$.
\end{claim}
\begin{proof}
Set $A=\arctan a$, $B=\arctan b$, $C=\arctan c$; then $A,B,C\in(0,\frac{\pi}{2})$. The condition $a+b+c=abc$ becomes $\tan A+\tan B+\tan C=\tan A\tan B\tan C$, which is equivalent to $A+B+C=\pi$. Moreover,
\[
\frac{1}{\sqrt{1+a^{2}}}=\cos A,\qquad
\frac{1}{\sqrt{1+b^{2}}}=\cos B,\qquad
\frac{1}{\sqrt{1+c^{2}}}=\cos C.
\]
The function $\cos x$ is concave on $[0,\frac{\pi}{2}]$ because its second derivative $-\cos x\le0$ there. By Jensen's inequality for concave functions,
\[
\frac{\cos A+\cos B+\cos C}{3}\le\cos\!\left(\frac{A+B+C}{3}\right)=\cos\frac{\pi}{3}=\frac12.
\]
Hence $\cos A+\cos B+\cos C\le\frac32$, i.e.
\[
\frac{1}{\sqrt{1+a^{2}}}+\frac{1}{\sqrt{1+b^{2}}}+\frac{1}{\sqrt{1+c^{2}}}\le\frac{3}{2}.
\]
Equality in Jensen holds when $A=B=C$. Together with $A+B+C=\pi$ this gives $A=B=C=\frac{\pi}{3}$, so $a=b=c=\tan\frac{\pi}{3}=\sqrt{3}$. Conversely, if $a=b=c=\sqrt{3}$, the left‑hand side equals $\frac{3}{\sqrt{1+3}}=\frac{3}{2}$.
\end{proof}

\end{document}